% Imporditud: tools/import_eksperimendid.py
% Allikas: Eesti füüsikaolümpiaadi eksperimendi ülesannete
%   kogu (Rael Kalda, 2023), ülesanne Ü164, lahendus L164.
\setAuthor{EFO žürii}
\setRound{piirkonnavoor}
\setYear{2005}
\setNumber{G E2}
\setDifficulty{5}
\setTopic{dünaamika}
\setVahendid{Nöör, horisontaalne laud, mõõtejoonlaud}
\setPunktid{16}
\setAllikas{Eksperimendi kogu Ü164, lahendus lk 122}
\setLahendusseis{toores}

\prob{Nöör}
Antud on homogeense massijaotusega nöör, mis asub horisontaalsel laual. Määrata nööri ja laua vaheline hõõrdetegur. Hinnata mõõtmisviga.

\hint

\solu
Lahendus 1 (vale): Selle lahenduse puhul me ignoreerime hõõret laua serva vastu, mis tegelikult on oluline. Katse idee (2 p). Olgu nööri pikkus l. Lasta nööri üks ots rippu üle laua serva, selle pikkus olgu $l_2$. Määrata laual lebava osa pikkus $l_1$ hetkel kui nöör hakkab laualt maha libisema.

Teooria: Nööri joontihedus (pikkusühiku mass) olgu $\rho$. Maksimaalne seisuhõõrdejõud Fh on võrdne rippuva nööriosa raskusjõuga

Fh = $\rho$gl2 (1 p).

Nööri rõhumisjõud lauale võrdub nööriosa raskusjõuga, mille pikkus on $l_1$:

Fh = $\rho$gl1 (1 p).

Seega hõõrdetegur $\mu$:

$\mu$ = Fh = $l_2$ (1 p). F $l_1$

Katse ja mõõtmised: Katsearuanne on vormistatud arusaadavalt (1 p). Katset on korratud mitu korda (1 p). Hinnatud on mõõteviga (2 p). Tulemus on tõepärane (1 p).

Numbriline näide: l = 45 cm, $l_1$ = 17 cm, $l_2$ = 28 cm. Seega $\mu$ = 1,65. Imeliku vastuse ($\mu$ > 1) põhjus on selles, et me ei arvestanud hõõret vastu nurka. Kokku ülesande eest maksimaalselt 10 p.

Lahendus 2:

Mõõtmised samad, mis lahenduses 1 (katse idee endiselt (2 p)), aga arvutame hoolikamalt. Panen tähele, et nöör ei paindu ümber nurga, vaid oma mõningase jäikuse tõttu puudutab laua nurka umbes 45-kraadise nurga all (horisondi suhtes) (1 p). Seega valem T $\propto$ e$\mu\varphi$ ei tööta (see eeldab, et nöör jälgib nurga kõverust, st liibub). Tuletame uue valemi lähtudes joonisest.

Esitame võrrandid, kus $\rho$g on välja taandatud.

2l2 = ($\mu$ + 1)N (2 p)

$\mu$$l_1$ + $l_2$ = 2N (2 p).

Siit leiame

$\mu$2l1 + $\mu$($l_2$ + $l_1$) $-$ $l_2$ = 0 (2 p).

Niisiis $\mu$ = l22 + l12 + 6l1l2 $-$ 2($l_2$ + $l_1$) (1 p). 2l1

Lahendusest 1 võetud mõõtmistulemuste järgi

$\mu \approx$ 0,52.

Katse ja mõõtmised: katsearuanne on vormistatud arusaadavalt (1 p). Katset on korratud mitu korda (1 p). Hinnatud on mõõteviga (2 p). Tulemus on tõepärane (2 p). Kokku ülesande eest maksimaalselt 16 p.

Lahendus 3: Lohistame nööri piki pinda ja teostame joonisel näidatud mõõtmised. Katse idee: 2 p.

Leiame: $l_1$ = 18 cm, rippuva osa pikkus $l_2$ = L $-$ $l_1$ = 27 cm, h = 20 cm, a = 7 cm. Paneme tähele, et pikkuse a mõõtmistulemus on üsna ebatäpne (kasutame nööri puutujana joonlauda). Paneme tähele, et rippuva nööri pinge horisontaalkomponendi konstantsusest johtuvalt

$\mu$gl1 = T cos $\alpha$ = gl2 cot $\alpha$ (3 p),

kus cot $\alpha$ = a/h (1 p). Niisiis $\mu$ = l2a/l1h $\approx$ 0,5. Katse ja mõõtmised: katsearuanne on vormistatud arusaadavalt (1 p). Katset on korratud mitu korda (1 p). Hinnatud on mõõteviga (2 p). Tulemus on tõepärane (2 p). Kokku ülesande eest maksimaalselt 12 p.

Märkus: madal punktisumma on motiveeritud antud meetodi madala täpsusega.
\probend
